\documentclass[10pt,journal]{IEEEtran}
\usepackage{graphicx}
\usepackage{amsmath,amssymb}
\usepackage{booktabs}
\usepackage{array}
\usepackage{hyperref}
\bibliographystyle{IEEEtaes.bst}

\renewcommand{\thefigure}{S\arabic{figure}}
\renewcommand{\thetable}{S\arabic{table}}
\renewcommand{\theequation}{S\arabic{equation}}
\renewcommand{\thesection}{S\arabic{section}}

\title{Supplementary Material for\\``Dual-Funnel Adaptive Sliding-Mode Control for Vision-Only Soft Precise Landing with Guaranteed Target Visibility''}
\author{Shubham Singhal,~Suresh Sundaram,~and~Jishnu Keshavan}

\begin{document}
\maketitle

\begin{abstract}
This document is the supplementary material accompanying the main manuscript. It contains (i) the virtual-camera abstraction and the baseline cascaded pose-based architecture whose inner-loop stages are reused by the proposed MDF-ASMC, (ii) detailed derivations including optic-flow funnel implementation remarks, SO(3) construction details, the yaw-channel Lyapunov substitution, and the full proof of Theorem~1, and (iii) the complete simulation result set: idealized-condition and realistic-disturbance figure sets, closed-loop internals, deep-sweep gain analysis, and the per-controller failure analysis of the comparative benchmark. Section, equation, figure, and table numbers in this document are prefixed with ``S''; references to the main paper use unprefixed numbers. All references to bibliography entries resolve against the main paper's bibliography.
\end{abstract}

% =================================================================
%  S1 — Virtual-Camera Model and Baseline (old S6)
% =================================================================
\section{Virtual-Camera Model and Baseline Cascaded Pose-Based Controller}
\label{sup:baseline}
This section collects the preliminary material that the main paper defers to the supplement: (i) the virtual-camera abstraction and the pinhole projection used throughout the image-based derivation, and (ii) the standard cascaded pose-based architecture whose inner-loop attitude and rate stages are reused verbatim by the proposed MDF-ASMC.

\subsection{Virtual-Camera Model and Feature Points}
\label{sup:virtual_features}
To decouple the translational image-feature dynamics from the rotational motion of the underactuated quadrotor, we adopt the virtual-camera abstraction of \cite{fink2017, xie2016_2}. Consider a camera frame $\mathcal{C}=\{O_{\text{c}}~X_{\text{c}}~Y_{\text{c}}~Z_{\text{c}}\}$ rigidly attached to the onboard monocular camera with its $X_{\text{c}}$ and $Y_{\text{c}}$ axes aligned with the image horizontal and vertical directions and $Z_{\text{c}}$ along the optical axis. Under the standard pinhole camera model \cite{chaumette2006}, a feature point on the image plane is
\begin{equation}
\label{sup:pinhole camera model}
    \,^\mathcal{C}\boldsymbol{\hat{r}} = \begin{bmatrix} \,^\mathcal{C}\hat{x} \\ \,^\mathcal{C}\hat{y} \end{bmatrix}
    = \frac{1}{\,^\mathcal{C}z}\begin{bmatrix} \,^\mathcal{C}x \\ \,^\mathcal{C}y \end{bmatrix},
\end{equation}
i.e.\ the perspective projection of the 3-D point $\,^\mathcal{C}\boldsymbol{r}=[\,^\mathcal{C}x,\,^\mathcal{C}y,\,^\mathcal{C}z]^{\top}$.

The virtual camera frame $\mathcal{V}=\{O_{\text{v}}~X_{\text{v}}~Y_{\text{v}}~Z_{\text{v}}\}$ moves with $\mathcal{C}$ but is constrained so that its $X_{\text{v}}$--$Y_{\text{v}}$ plane remains parallel to the inertial $X_{\text{i}}$--$Y_{\text{i}}$ plane; the associated virtual image plane keeps the same position and orientation relative to $\mathcal{V}$ as the physical image plane has relative to $\mathcal{C}$. The corresponding rotation is obtained by rolling out $\phi$ and pitching out $\theta$:
\begin{equation}
\label{sup:virtual 3D point}
    \,^\mathcal{V}\boldsymbol{r} = \,^\mathcal{V}R_\mathcal{C}\,^\mathcal{C}\boldsymbol{r}, \qquad \,^\mathcal{V}R_\mathcal{C} = R_y(\theta)R_x(\phi).
\end{equation}
Combining \eqref{sup:virtual 3D point} with the pinhole model \eqref{sup:pinhole camera model}, the virtual feature point is
\begin{equation}
\label{sup:virtual feature point}
    \,^\mathcal{V}\boldsymbol{\hat{r}} = \begin{bmatrix} \,^\mathcal{V}\hat{x} \\ \,^\mathcal{V}\hat{y} \end{bmatrix}
    = \frac{1}{\,^\mathcal{V}z}\begin{bmatrix} \,^\mathcal{V}x \\ \,^\mathcal{V}y \end{bmatrix}
    = \begin{bmatrix}
        \dfrac{C_\theta\,^\mathcal{C}\hat{x} + S_\theta S_\phi\,^\mathcal{C}\hat{y} + S_\theta C_\phi}{-S_\theta\,^\mathcal{C}\hat{x} + C_\theta S_\phi\,^\mathcal{C}\hat{y} + C_\theta C_\phi} \\[4pt]
        \dfrac{C_\theta\,^\mathcal{C}\hat{y} - S_\phi}{-S_\theta\,^\mathcal{C}\hat{x} + C_\theta S_\phi\,^\mathcal{C}\hat{y} + C_\theta C_\phi}
    \end{bmatrix},
\end{equation}
with $S_x=\sin(x)$, $C_x=\cos(x)$. Throughout the main paper, $\,^\mathcal{V}\boldsymbol{\hat{r}}$ and the normalized position $\boldsymbol{s}=[\,^\mathcal{V}\boldsymbol{\hat{r}}^\top,\,1]^\top$ are the scale-independent image features regulated by MDF-ASMC.

\subsection{Cascaded Pose-Based Architecture}
\label{sup:cascaded_baseline}
The proposed image-based flight controller of the main paper reuses the inner-loop attitude and rate stages of the standard cascaded pose-based architecture, but replaces its outer-loop pose feedback with image-feature feedback. For completeness we recapitulate the baseline architecture here.

In order to achieve target tracking using a pose-based flight controller, the overall control architecture is a cascaded PID controller with the outer loop operating in the inertial frame $\mathcal{I}=\{O_{\text{i}}~X_{\text{i}}~Y_{\text{i}}~Z_{\text{i}}\}$ and the inner loop in the body frame $\mathcal{B}$. The outer-loop controller generates a desired velocity $\,^\mathcal{I}\boldsymbol{v}_{\text{d}}$ command that drives the UAV from its instantaneous position $\,^\mathcal{I}\boldsymbol{r}_{\text{b}}$ to the known target position $\,^\mathcal{I}\boldsymbol{r}_{\text{t}}$, followed by a desired acceleration $\,^\mathcal{I}\boldsymbol{a}_{\text{d}}$ command that drives the instantaneous UAV velocity $\,^\mathcal{I}\boldsymbol{v}_{\text{b}}$ to $\,^\mathcal{I}\boldsymbol{v}_{\text{d}}$:
\begin{align} \label{sup:outer_loop_control}
    \,^\mathcal{I}\boldsymbol{v}_{\text{d}} =& - K_{rp} \,^\mathcal{I}\boldsymbol{r}_{\text{e}} - K_{ri} \int \,^\mathcal{I}\boldsymbol{r}_{\text{e}}\, dt - K_{rd} \frac{d\,^\mathcal{I}\boldsymbol{r}_{\text{e}}}{dt}, \nonumber \\
    \,^\mathcal{I}\boldsymbol{a}_{\text{d}} =&  K_{vp} \,^\mathcal{I}\boldsymbol{v}_{\text{e}} - K_{vi} \int \,^\mathcal{I}\boldsymbol{v}_{\text{e}}\, dt - K_{vd} \frac{d\,^\mathcal{I}\boldsymbol{v}_{\text{e}}}{dt},
\end{align}
where $\,^\mathcal{I}\boldsymbol{r}_{\text{e}} = \,^\mathcal{I}\boldsymbol{r}_{\text{b}} - \,^\mathcal{I}\boldsymbol{r}_{\text{t}}$ is the relative position between UAV and target and $\,^\mathcal{I}\boldsymbol{v}_{\text{e}} = \,^\mathcal{I}\boldsymbol{v}_{\text{b}} - \,^\mathcal{I}\boldsymbol{v}_{\text{d}}$.

Since the centrifugal force $\,^\mathcal{B}\boldsymbol{\omega}_{\text{b}}\times(m\,^\mathcal{B}\boldsymbol{v}_{\text{b}})$ is nullified in the inertial frame, the desired acceleration $\,^\mathcal{I}\boldsymbol{a}_{\text{d}}$ is assumed to be achieved solely by the total rotor thrust $\,^\mathcal{B}T_u$ and its direction after compensating for gravity and disturbances:
\begin{align} \label{sup:a_des}
    \,^\mathcal{I}\boldsymbol{a}_{\text{d}} &= \frac{1}{m}\,^\mathcal{I}R_\mathcal{B}\,^\mathcal{B}\boldsymbol{F}_u = \frac{1}{m}\,^\mathcal{I}R_\mathcal{B}\begin{bmatrix} 0 & 0 & \,^\mathcal{B}T_u \end{bmatrix}^\top,
\end{align}
where $\,^\mathcal{I}R_\mathcal{B}$ is the rotation matrix from $\mathcal{B}$ to $\mathcal{I}$. Inverting \eqref{sup:a_des} using the components $\,^\mathcal{I}a_{\text{d}_x},\,^\mathcal{I}a_{\text{d}_y},\,^\mathcal{I}a_{\text{d}_z}$ yields the desired roll, pitch, and total-thrust commands:
\begin{align} \label{sup:euler_thrust_des}
    \theta_{\text{d}} &= \tan^{-1}\!\left(\tfrac{-S_\psi\,^\mathcal{I}a_{\text{d}_x} - C_\psi\,^\mathcal{I}a_{\text{d}_y}}{\,^\mathcal{I}a_{\text{d}_z}}\right), \nonumber \\
    \phi_{\text{d}} &= \tan^{-1}\!\left(\tfrac{C_{\theta_\text{d}}(C_\psi\,^\mathcal{I}a_{\text{d}_x} - S_\psi\,^\mathcal{I}a_{\text{d}_y})}{\,^\mathcal{I}a_{\text{d}_z}}\right), \nonumber \\
    \,^\mathcal{B}T_u &= -\tfrac{m\,^\mathcal{I}a_{\text{d}_z}}{C_{\theta_\text{d}}C_{\phi_\text{d}}},
\end{align}
where $S_x=\sin(x)$, $C_x=\cos(x)$. Combined with the known target heading $\psi_\text{t}$, the desired attitude $\boldsymbol{\Phi}_\text{d}=[\phi_\text{d},\theta_\text{d},\psi_\text{t}]^\top$ and thrust $\,^\mathcal{B}T_u$ are passed to the inner-loop controller.

The inner-loop controller regulates the attitude $\boldsymbol{\Phi}_\text{b}$ and body rates $\,^\mathcal{B}\boldsymbol{\omega}_\text{b}$ through a cascade of PID loops. The desired attitude rates $\dot{\boldsymbol{\Phi}}_\text{d}$ are mapped to desired body rates $\,^\mathcal{B}\boldsymbol{\omega}_\text{d}$ via the kinematic transformation
\begin{equation} \label{sup:kinematic_transformation}
    \,^\mathcal{B}\boldsymbol{\omega}_\text{d} = W(\phi_\text{d},\theta_\text{d})\dot{\boldsymbol{\Phi}}_\text{d},
\end{equation}
with
\begin{equation*}
    W(\phi_\text{d},\theta_\text{d}) = \begin{bmatrix} 0 & C_{\phi_\text{d}} & S_{\phi_\text{d}}C_{\theta_\text{d}} \\ 1 & 0 & -S_{\theta_\text{d}} \\ 0 & S_{\phi_\text{d}} & -C_{\phi_\text{d}}C_{\theta_\text{d}} \end{bmatrix}.
\end{equation*}
The inner-loop PID control laws are
\begin{align}
    \dot{\boldsymbol{\Phi}}_\text{d} &= - K_{\Phi p}\boldsymbol{\Phi}_\text{e} - K_{\Phi i}\!\int\!\boldsymbol{\Phi}_\text{e}\,dt - K_{\Phi d}\frac{d\boldsymbol{\Phi}_\text{e}}{dt}, \nonumber \\
    \,^\mathcal{B}\dot{\boldsymbol{\omega}}_\text{d} &= - K_{\omega p}\,^\mathcal{B}\boldsymbol{\omega}_\text{e} - K_{\omega i}\!\int\!^\mathcal{B}\boldsymbol{\omega}_\text{e}\,dt - K_{\omega d}\frac{d\,^\mathcal{B}\boldsymbol{\omega}_\text{e}}{dt},
\end{align}
with $\boldsymbol{\Phi}_\text{e}=\boldsymbol{\Phi}_\text{b}-\boldsymbol{\Phi}_\text{d}$ and $\,^\mathcal{B}\boldsymbol{\omega}_\text{e}=\,^\mathcal{B}\boldsymbol{\omega}_\text{b}-\,^\mathcal{B}\boldsymbol{\omega}_\text{d}$. The desired control torque is finally recovered from the rigid-body dynamics as
\begin{equation} \label{sup:tau_des}
    \,^\mathcal{B}\boldsymbol{\tau}_u = J\,^\mathcal{B}\dot{\boldsymbol{\omega}}_\text{d} + \,^\mathcal{B}\boldsymbol{\omega}_\text{d}\times(J\,^\mathcal{B}\boldsymbol{\omega}_\text{d}),
\end{equation}
and the final motor-mixer input is the combined thrust-torque command $\boldsymbol{u}=[\,^\mathcal{B}T_u,\,^\mathcal{B}\boldsymbol{\tau}_u^\top]^\top$. The proposed controller in the main paper inherits \eqref{sup:kinematic_transformation}--\eqref{sup:tau_des} verbatim, replacing only the heading $\psi$ in the inner cascade with the image-plane orientation $\alpha$.

% =================================================================
%  S2 — Detailed Derivations of MDF-ASMC (old S7)
% =================================================================
\clearpage
\section{Detailed Derivations of MDF-ASMC}
\label{sup:detailed}
This section collects the standard prescribed-performance machinery, implementation remarks, SO(3) construction details, the yaw-channel Lyapunov substitution, and the full proof of Theorem~1 that the main paper defers to the supplement.

\subsection{Optic-Flow Funnel: Implementation Remarks}
\label{sup:optic_flow_remarks}
\textit{Remark (Initial-condition feasibility).} The optic-flow transformation requires $|h_{e_k}(0)|<p_{2_{0k}}$. In the reported runs, $\boldsymbol{p}_{2_0}=[25,25,8]^\top$ is sized to envelope the worst-case initial optic-flow magnitude observed across the five ICs of Section~IV-A of the main paper, with a $2\times$--$3\times$ margin against pixel-noise spikes on the fastest \texttt{traj\_Gen} targets (Case~5 at $\omega_z=0.3$~rad/s; Case~2 at $1.414$~m/s).

\textit{Remark (Sliding-surface integral anti-windup).} At the funnel mouth, $\int\boldsymbol{\zeta}_2\,d\tau$ in the sliding-surface definition can wind up whenever $\boldsymbol{\zeta}_2$ is briefly clamped by the saturation guard of Remark~1 in the main paper. We therefore bound each channel by $|\!\int\!\zeta_{2_k}\,d\tau|\le 10$. The clamp activates only during short transient spikes and does not affect the Lyapunov derivative on the open sub-funnel where Theorem~1 operates.

\subsection{Yaw-Channel Lyapunov Substitution}
\label{sup:yaw_lyap}
Because the yaw channel decouples from the translational dynamics under the virtual-camera abstraction, the Lyapunov argument of Theorem~1 applies to the scalar yaw ASMC with the substitutions $\boldsymbol{\sigma}\leftarrow\sigma_a$, $\Gamma\leftarrow\gamma_a$, $\mathcal{P}\leftarrow p_a$, $\mathcal{N}\leftarrow n_a$, $\mathcal{G}_2\leftarrow 1$, and $\beta_a\equiv 1$ (unity heading-rate gain under the virtual-compass integrator). The scalar Lyapunov function $V_a(t) = \tfrac12\sigma_a^2 + \tfrac12 n_a^{-1}(\kappa_a-\tilde{d}_a)^2$ reproduces the main-paper derivative inequalities line for line with $\lambda_{\min}(\Gamma)\leftarrow\gamma_a$, $\lambda_{\max}(\mathcal{P})\leftarrow p_a$, and yields the yaw-channel UUB
\begin{equation*}
    \vartheta_a = \sqrt{\tfrac{3 p_a}{\varphi_{1a}-\varphi_{2a}}}\,\tilde{d}_a,
\end{equation*}
which renders $\alpha_e$ uniformly ultimately bounded. The virtual-compass integrator $\psi_\text{d}(t_{k+1}) = \Pi_{[-\pi,\pi]}[\psi_\text{d}(t_k)+\dot{\psi}_\text{d}(t_k)\Delta t]$ then yields a bounded desired heading $\psi_\text{d}$ for the inner SO(3) tracker. The raw image-feature angle is unwrapped to $(-\pi/2,\pi/2]$ via $\alpha_e = \tfrac12\text{atan2}(\sin(2\alpha_e^\text{raw}),\cos(2\alpha_e^\text{raw}))$ to absorb the ellipse-orientation half-turn symmetry.

\subsection{SO(3) Inner-Loop Construction Details}
\label{sup:so3_details}
The outer adaptive law returns a desired inertial acceleration $\,^\mathcal{I}\!\boldsymbol{a}_\text{d}$ that is first low-pass filtered to absorb pixel-noise spikes amplified through $L_s^{-1}$ at low altitude:
\begin{equation} \label{sup:I a cd filt}
    \,^\mathcal{I}\!\tilde{\boldsymbol{a}}_\text{d}(t_{k+1}) = \alpha_\text{a}\,^\mathcal{I}\!\tilde{\boldsymbol{a}}_\text{d}(t_k)+(1-\alpha_\text{a})\,^\mathcal{I}\!\boldsymbol{a}_\text{d}(t_{k+1}),
\end{equation}
with $\alpha_\text{a}=\tau_\text{a}/(\tau_\text{a}+\Delta t)$ and $\tau_\text{a}=0.08$~s ($\sim 2$~Hz cut-off). The commanded thrust magnitude and desired body-$z$ axis are then $\,^\mathcal{B}T_u = m\|\,^\mathcal{I}\!\tilde{\boldsymbol{a}}_\text{d}\|$ and $\boldsymbol{r}_{d,3} = -\,^\mathcal{I}\!\tilde{\boldsymbol{a}}_\text{d}/\|\,^\mathcal{I}\!\tilde{\boldsymbol{a}}_\text{d}\|$. With $\psi_\text{d}$ from the virtual-compass integrator supplying $\boldsymbol{a}_h=[\cos\psi_\text{d},\sin\psi_\text{d},0]^\top$, the remaining body axes are constructed by Gram--Schmidt on $SO(3)$:
\begin{equation} \label{sup:R_d construction}
    \boldsymbol{r}_{d,2} = \frac{\boldsymbol{r}_{d,3}\times\boldsymbol{a}_h}{\|\boldsymbol{r}_{d,3}\times\boldsymbol{a}_h\|},\ \ \boldsymbol{r}_{d,1} = \boldsymbol{r}_{d,2}\times\boldsymbol{r}_{d,3},\ \ R_\text{d}=[\boldsymbol{r}_{d,1}\ \boldsymbol{r}_{d,2}\ \boldsymbol{r}_{d,3}].
\end{equation}
A degeneracy guard replaces $\boldsymbol{r}_{d,3}\times\boldsymbol{a}_h$ with a fixed east-axis fallback when $\|\boldsymbol{r}_{d,3}\times\boldsymbol{a}_h\|<10^{-6}$. The attitude and angular-velocity errors are
\begin{equation*}
    \boldsymbol{e}_R = \tfrac12(R_\text{d}^\top R - R^\top R_\text{d})^\vee,\qquad \boldsymbol{e}_\Omega = \boldsymbol{\omega} - R^\top R_\text{d}\boldsymbol{\omega}_\text{d},
\end{equation*}
with $\boldsymbol{\omega}_\text{d}\equiv 0$ since the outer loop commands attitude (not attitude rate), so $\boldsymbol{e}_\Omega=\boldsymbol{\omega}$. The geometric torque law of the main paper then inherits the almost-global exponential stability result of \cite[Thm.~3.1]{lee2010}: for any initial rotation outside the co-dimension-1 set $\{R:\text{tr}(R_\text{d}^\top R)=-1\}$, $(R_\text{d}^\top R,\boldsymbol{e}_\Omega)$ decays exponentially to $(I,\boldsymbol{0})$.

\subsection{Full Proof of Theorem~1}
\label{sup:theorem1_proof}
\textit{Definition (GUUB).} A multivariable signal $\boldsymbol{\chi}(t)\in\mathbb{R}^n$ is \emph{globally uniformly ultimately bounded} by $\Lambda>0$ if for every $\lambda>0$ there exists $\mathcal{T}\ge 0$ such that $\|\boldsymbol{\chi}(0)\|\le\lambda\Rightarrow\|\boldsymbol{\chi}(t)\|\le\Lambda\ \forall t\ge\mathcal{T}$.

\textit{Lemma.} If $V(\boldsymbol{\chi})$ is continuously differentiable with $k_1\|\boldsymbol{\chi}\|^2\le V\le k_2\|\boldsymbol{\chi}\|^2$ and $\dot{V}\le -k_3 V$ for $\|\boldsymbol{\chi}\|\ge\varkappa$ ($k_1,k_2,k_3>0$), then every trajectory with $\|\boldsymbol{\chi}(0)\|^2\le\lambda$ is GUUB with $\Lambda = \varkappa(k_2/k_1)^{1/2}$.

\textit{Proof of Theorem~1.} The initial transformed error $\zeta_{2_k}(0) = S_{2_k}^{-1}(h_{e_k}(0)/p_{2_{0k}})$ is well defined since $|h_{e_k}(0)|<p_{2_{0k}}$. Under the control law of the main paper, $\boldsymbol{\zeta}_2$ evolves via the unconstrained dynamics equation of the main paper (Section~III-A2). Consider the candidate
\begin{equation} \label{sup:Lyapunov function}
    V(t) = \tfrac12\boldsymbol{\sigma}^\top\boldsymbol{\sigma} + \tfrac{\beta_{\min}}{2}[\boldsymbol{\kappa}(t)-\tilde{d}\boldsymbol{1}]^\top\mathcal{N}^{-1}[\boldsymbol{\kappa}(t)-\tilde{d}\boldsymbol{1}],
\end{equation}
with $\boldsymbol{1}=[1,1,1]^\top$ and $\|\boldsymbol{\overline{d}}\|\le\tilde{d}$. Using $\dot{\boldsymbol{\sigma}}$ from the main paper and the adaptive law,
\begin{align}
\label{sup:Vdot 1}
\dot{V}(t) &= \boldsymbol{\sigma}^\top\dot{\boldsymbol{\sigma}} + \beta_{\min}[\boldsymbol{\kappa}(t)-\tilde{d}\boldsymbol{1}]^\top\mathcal{N}^{-1}\dot{\boldsymbol{\kappa}}(t)\nonumber\\
&= -\beta\,\boldsymbol{\sigma}^\top\!\big[\Gamma\boldsymbol{\sigma}+\|\boldsymbol{\theta}\|\text{sat}(\mathcal{E}^{-1}\boldsymbol{\sigma})\mathcal{G}_2\boldsymbol{\kappa}(t)\big]\nonumber\\
&\quad + \beta_{\min}\boldsymbol{\sigma}^\top\mathcal{G}_2\boldsymbol{\theta}\boldsymbol{\overline{d}}\nonumber\\
&\quad + \beta_{\min}[\boldsymbol{\kappa}(t)-\tilde{d}\boldsymbol{1}]^\top\mathcal{N}^{-1}\!\big[\|\boldsymbol{\theta}\|\mathcal{N}\mathcal{G}_2|\boldsymbol{\sigma}| - \mathcal{N}\mathcal{P}\boldsymbol{\kappa}(t)\big].
\end{align}
Since $\sigma_k\text{sat}(\sigma_k/\varepsilon_k)\ge 0$ componentwise and $\boldsymbol{\kappa}(t)\in\mathbb{R}_+^3$, the switching dissipation satisfies
\begin{equation*}
\boldsymbol{\sigma}^\top\text{sat}(\mathcal{E}^{-1}\boldsymbol{\sigma}) = |\boldsymbol{\sigma}|^\top - \boldsymbol{r}_\varepsilon^\top(\boldsymbol{\sigma}),
\end{equation*}
where $r_{\varepsilon_k} = \max(0,|\sigma_k|-\sigma_k^2/\varepsilon_k)\le\varepsilon_k/4$ vanishes outside the boundary layer $|\sigma_k|>\varepsilon_k$. The residual $\tfrac14\|\boldsymbol{\varepsilon}\|\lambda_{\max}(\mathcal{G}_2)\|\boldsymbol{\kappa}(t)\|$ is bounded uniformly by the leakage law and absorbed into $\tilde{d}$. Using $\beta\ge\beta_{\min}>0$, \eqref{sup:Vdot 1} rearranges to
\begin{align*}
\dot{V}(t)\le& -\beta_{\min}\boldsymbol{\sigma}^\top\Gamma\boldsymbol{\sigma} - \beta_{\min}(\tilde{d}-\|\boldsymbol{\overline{d}}\|)\|\boldsymbol{\theta}\||\boldsymbol{\sigma}|^\top\mathcal{G}_2\boldsymbol{1}\\
&- \beta_{\min}[\boldsymbol{\kappa}^\top(t)\mathcal{P}\boldsymbol{\kappa}(t) - \tilde{d}\boldsymbol{1}^\top\mathcal{P}\boldsymbol{\kappa}(t)].
\end{align*}
The second term is non-negative since $\|\boldsymbol{\overline{d}}\|\le\tilde{d}$. Completing the square on the third term and using positive-definiteness of $\mathcal{P}$,
\begin{align*}
\dot{V}(t)&\le -\beta_{\min}\lambda_{\min}(\Gamma)\|\boldsymbol{\sigma}\|^2 - \tfrac{\beta_{\min}}{2}[\boldsymbol{\kappa}(t)-\tilde{d}\boldsymbol{1}]^\top\mathcal{P}[\boldsymbol{\kappa}(t)-\tilde{d}\boldsymbol{1}]\\
&\quad + \tfrac{\beta_{\min}}{2}\tilde{d}^2\boldsymbol{1}^\top\mathcal{P}\boldsymbol{1}\\
&\le -\varphi_1 V(t) + \tfrac32\beta_{\min}\lambda_{\max}(\mathcal{P})\tilde{d}^2.
\end{align*}
For any $0<\varphi_2<\varphi_1$,
\begin{equation*}
\dot{V}(t)\le -\varphi_2 V(t),\qquad \forall\,V(t)\ge\Theta=\tfrac{3\beta_{\min}\lambda_{\max}(\mathcal{P})\tilde{d}^2}{2(\varphi_1-\varphi_2)}.
\end{equation*}
Invoking the Lemma, every trajectory enters the ball $\{V\le\Theta\}$ in finite time and remains there. Since $V(t)\ge\|\boldsymbol{\sigma}\|^2/2$, the ultimate bound on $\boldsymbol{\sigma}$ is $\vartheta = [3\beta_{\min}\lambda_{\max}(\mathcal{P})/(\varphi_1-\varphi_2)]^{1/2}\tilde{d}$, independent of initial conditions. From the sliding surface definition, $\boldsymbol{\zeta}_2$ converges into a small residual around the origin, so from the inverse PPC transformation $-\boldsymbol{p}_2(t)<\boldsymbol{h}_e<\boldsymbol{p}_2(t)$ holds for all $t\ge 0$. $\hfill\blacksquare$

\textit{Cascade argument for Corollary~1.} Theorem~1 renders $\boldsymbol{h}_\text{e}=\boldsymbol{h}-\boldsymbol{h}_\text{d}$ uniformly ultimately bounded inside $\boldsymbol{p}_2(t)$. Combining the visual error kinematics of the main paper with the PID-generated reference and substituting $\boldsymbol{h}=\boldsymbol{h}_\text{d}+\boldsymbol{h}_\text{e}$ yields the horizontal-channel closed loop
\begin{equation*}
    (P_{1_0}+K_{rd})\ddot{\bar{\boldsymbol{r}}}_\text{e} + K_{rp}\dot{\bar{\boldsymbol{r}}}_\text{e} + K_{ri}\bar{\boldsymbol{r}}_\text{e} = \boldsymbol{\delta}_\text{e}(\boldsymbol{h}_\text{e},\boldsymbol{s}),
\end{equation*}
with $P_{1_0}=\text{diag}(\boldsymbol{p}_{1_0})$ and the perturbation $\boldsymbol{\delta}_\text{e}$ linear in $\boldsymbol{h}_\text{e}$ and bounded for all $\boldsymbol{s}$ inside the FoV. This closed loop is a linear second-order PID-driven system; under the diagonal positive-definite gains $(K_{rp},K_{ri},K_{rd})$ every channel is Hurwitz. Standard ISS arguments for linear cascades~\cite{khalil2002} therefore render $\bar{\boldsymbol{r}}_\text{e}$, and equivalently $\hat{\boldsymbol{r}}_\text{e}=\boldsymbol{p}_{1_0}\odot\bar{\boldsymbol{r}}_\text{e}$, ultimately bounded by an affine function of $\vartheta$. The visibility constraint is enforced kinematically by the FoV-adaptive cone clamp of Section~IV of the main paper, which guarantees that every target image corner point stays inside $\boldsymbol{\rho}_\text{fov}(t)\subseteq\boldsymbol{\rho}_\text{fov}(0)\preceq\mathcal{R}/2$ for all $t\ge 0$ independently of the outer-loop transient.

% =================================================================
%  S3 — Simulation Results (old S1–S5 merged)
% =================================================================
\clearpage
\section{Simulation Results}
\label{sec:sup-results}
This section collects the complete simulation result set supporting Section~IV of the main paper: the idealized-condition and realistic-disturbance multi-initial-condition figure sets, the closed-loop internals of a representative run, the deep-sweep gain analysis, and the per-controller failure analysis of the comparative benchmark.

\subsection{Idealized-Condition Multi-Initial-Condition Trajectories}
\label{sec:sup-noiseless}

The figures in this section verify that, in the absence of measurement noise, ground effect, computational delay, and model uncertainty, the closed loop lands from every initial condition of Table~II of the main paper on every trajectory of Section~IV-A of the main paper, with the target image corner points remaining strictly inside the FoV-adaptive pixel-box $\boldsymbol{\rho}_\text{fov}(t)$. These runs provide an empirical verification of the dual-funnel construction, isolating the role of the prescribed-performance layer from that of the adaptive robustness layer.

\begin{figure}[!t]
\centering
    \includegraphics[width=\columnwidth]{Figures/generated/multi_init/Static_3D_noiseless.pdf}
    \caption{3-D UAV trajectories for Case~1 (static target) under idealized conditions.}
    \label{fig:sup-static-3d-nl}
\end{figure}

\begin{figure}[!t]
\centering
    \includegraphics[width=\columnwidth]{Figures/generated/multi_init/Static_image_plane_noiseless.pdf}
    \caption{Image-plane evolution for Case~1 (static target) under idealized conditions.}
    \label{fig:sup-static-ip-nl}
\end{figure}

\begin{figure}[!t]
\centering
    \includegraphics[width=\columnwidth]{Figures/generated/multi_init/Linear_3D_noiseless.pdf}
    \caption{3-D UAV trajectories for Case~2 (linear target trajectory) under idealized conditions.}
    \label{fig:sup-linear-3d-nl}
\end{figure}

\begin{figure}[!t]
\centering
    \includegraphics[width=\columnwidth]{Figures/generated/multi_init/Linear_image_plane_noiseless.pdf}
    \caption{Image-plane evolution for Case~2 (linear target trajectory) under idealized conditions.}
    \label{fig:sup-linear-ip-nl}
\end{figure}

\begin{figure}[!t]
\centering
    \includegraphics[width=\columnwidth]{Figures/generated/multi_init/Sinusoidal_3D_noiseless.pdf}
    \caption{3-D UAV trajectories for Case~3 (sinusoidal target trajectory) under idealized conditions.}
    \label{fig:sup-sin-3d-nl}
\end{figure}

\begin{figure}[!t]
\centering
    \includegraphics[width=\columnwidth]{Figures/generated/multi_init/Sinusoidal_image_plane_noiseless.pdf}
    \caption{Image-plane evolution for Case~3 (sinusoidal target trajectory) under idealized conditions.}
    \label{fig:sup-sin-ip-nl}
\end{figure}

\begin{figure}[!t]
\centering
    \includegraphics[width=\columnwidth]{Figures/generated/multi_init/Lissajous_3D_noiseless.pdf}
    \caption{3-D UAV trajectories for Case~4 (Lissajous target trajectory) under idealized conditions.}
    \label{fig:sup-liss-3d-nl}
\end{figure}

\begin{figure}[!t]
\centering
    \includegraphics[width=\columnwidth]{Figures/generated/multi_init/Lissajous_image_plane_noiseless.pdf}
    \caption{Image-plane evolution for Case~4 (Lissajous target trajectory) under idealized conditions.}
    \label{fig:sup-liss-ip-nl}
\end{figure}

\begin{figure}[!t]
\centering
    \includegraphics[width=\columnwidth]{Figures/generated/multi_init/Circular_3D_noiseless.pdf}
    \caption{3-D UAV trajectories for Case~5 (circular target trajectory) under idealized conditions.}
    \label{fig:sup-circ-3d-nl}
\end{figure}

\begin{figure}[!t]
\centering
    \includegraphics[width=\columnwidth]{Figures/generated/multi_init/Circular_image_plane_noiseless.pdf}
    \caption{Image-plane evolution for Case~5 (circular target trajectory) under idealized conditions.}
    \label{fig:sup-circ-ip-nl}
\end{figure}

\clearpage
\subsection{Full Multi-Initial-Condition Figure Set under Realistic Disturbances}
\label{sec:sup-multi-init}

The main paper shows a single representative Circular ship-deck run under realistic disturbances; the figures below provide the full per-trajectory set that the Table~V multi-init statistics of the main paper are aggregated over. The runs share the full disturbance model of Section~IV-B of the main paper: depth-dependent pixel noise $\sigma_{\text{px}}(z)=0.3+0.5/(z+0.5)$~px, $0.5\%$ Bernoulli $5$-px pixel outliers, mean wind $\pm 0.2$~m/s plus Ornstein--Uhlenbeck turbulence ($\tau=1$~s, $\sigma_w=0.1$~N), $\pm 5\%$ parametric uncertainty on mass and inertia, $\pm 5$~mm centre-of-gravity offset, aerodynamic drag $C_d=0.25$, rotor ground effect, one-step computational delay, and $35^\circ$ attitude-cone projection. Cases~2 and~5 additionally carry $0.2\sin(0.5\,t)~\text{m}$ heave and $(\pm 15^\circ,\pm 8^\circ)$ roll/pitch oscillations, matching the ship-deck profile of~\cite{lin2022}.

\begin{figure}[!t]
\centering
    \includegraphics[width=\columnwidth]{Figures/generated/multi_init/Static_3D.pdf}
    \caption{3-D UAV trajectories for Case~1 (static target) under realistic disturbances.}
    \label{fig:sup-static-3d}
\end{figure}

\begin{figure}[!t]
\centering
    \includegraphics[width=\columnwidth]{Figures/generated/multi_init/Static_image_plane.pdf}
    \caption{Image-plane evolution for Case~1 (static target) under realistic disturbances.}
    \label{fig:sup-static-ip}
\end{figure}

\begin{figure}[!t]
\centering
    \includegraphics[width=\columnwidth]{Figures/generated/multi_init/Linear_3D.pdf}
    \caption{3-D UAV trajectories for Case~2 (linear target trajectory) under realistic disturbances.}
    \label{fig:sup-linear-3d}
\end{figure}

\begin{figure}[!t]
\centering
    \includegraphics[width=\columnwidth]{Figures/generated/multi_init/Linear_image_plane.pdf}
    \caption{Image-plane evolution for Case~2 (linear target trajectory) under realistic disturbances.}
    \label{fig:sup-linear-ip}
\end{figure}

\begin{figure}[!t]
\centering
    \includegraphics[width=\columnwidth]{Figures/generated/multi_init/Sinusoidal_3D.pdf}
    \caption{3-D UAV trajectories for Case~3 (sinusoidal target trajectory) under realistic disturbances.}
    \label{fig:sup-sin-3d}
\end{figure}

\begin{figure}[!t]
\centering
    \includegraphics[width=\columnwidth]{Figures/generated/multi_init/Sinusoidal_image_plane.pdf}
    \caption{Image-plane evolution for Case~3 (sinusoidal target trajectory) under realistic disturbances.}
    \label{fig:sup-sin-ip}
\end{figure}

\begin{figure}[!t]
\centering
    \includegraphics[width=\columnwidth]{Figures/generated/multi_init/Lissajous_3D.pdf}
    \caption{3-D UAV trajectories for Case~4 (Lissajous target trajectory) under realistic disturbances.}
    \label{fig:sup-liss-3d}
\end{figure}

\begin{figure}[!t]
\centering
    \includegraphics[width=\columnwidth]{Figures/generated/multi_init/Lissajous_image_plane.pdf}
    \caption{Image-plane evolution for Case~4 (Lissajous target trajectory) under realistic disturbances.}
    \label{fig:sup-liss-ip}
\end{figure}

\begin{figure}[!t]
\centering
    \includegraphics[width=\columnwidth]{Figures/generated/multi_init/Circular_image_plane.pdf}
    \caption{Image-plane evolution for Case~5 (circular target trajectory) under realistic disturbances (the corresponding 3-D trajectory view appears as Fig.~9 in the main paper).}
    \label{fig:sup-circ-ip}
\end{figure}

\clearpage
\subsection{Closed-Loop Internals of MDF-ASMC (Representative Run)}
\label{sec:sup-internals}

The figures in this section expose the closed-loop internals of MDF-ASMC for the representative case IC$2=[2,2,-5]$~m on Case~3 (sinusoidal trajectory) under the full disturbance set. They provide the per-signal evidence underlying the single-sentence empirical claim in Section~IV-B of the main paper, namely that the outer image-feature error, the inner optic-flow error, the sliding surface, the adaptive gain, and the commanded thrust/acceleration all remain inside their respective safety envelopes at every time instant.

\begin{figure}[!t]
\centering
    \includegraphics[width=\columnwidth]{Figures/generated/plasmc_outer_funnel.pdf}
    \caption{Visibility funnel: image-feature centroid error $\boldsymbol{E}_2(t)$ versus the prescribed-performance bound $\pm\boldsymbol{p}_1(t)$.}
    \label{fig:sup-outer}
\end{figure}

\begin{figure}[!t]
\centering
    \includegraphics[width=\columnwidth]{Figures/generated/plasmc_inner_funnel.pdf}
    \caption{Optic-flow funnel: optic-flow tracking error $\boldsymbol{w}_e(t)$ versus $\pm\boldsymbol{p}_2(t)$ on all three axes.}
    \label{fig:sup-inner}
\end{figure}

\begin{figure}[!t]
\centering
    \includegraphics[width=\columnwidth]{Figures/generated/plasmc_sliding.pdf}
    \caption{Sliding-surface evolution $\boldsymbol{\sigma}(t)$ on each translational axis.}
    \label{fig:sup-sliding}
\end{figure}

\begin{figure}[!t]
\centering
    \includegraphics[width=\columnwidth]{Figures/generated/plasmc_adaptive_gain.pdf}
    \caption{Leakage-type adaptive gains $\boldsymbol{\kappa}(t)$ and yaw-channel gain $\kappa_\text{a}(t)$.}
    \label{fig:sup-kappa}
\end{figure}

\begin{figure}[!t]
\centering
    \includegraphics[width=\columnwidth]{Figures/generated/plasmc_thrust_accel.pdf}
    \caption{Total thrust command $\,^\mathcal{B}T_u(t)$ and lateral commanded acceleration $\|\,^\mathcal{I}\boldsymbol{a}_{\text{d},xy}\|$ versus the cone-induced limit.}
    \label{fig:sup-thrust}
\end{figure}

\clearpage
\subsection{Locked MDF-ASMC Control Parameters}
\label{sec:sup-gains}

Table~\ref{sup:control params} lists the complete locked MDF-ASMC gain set used in every numerical test reported in the main paper. These gains were obtained via the cross-trajectory deep sweep of Section~\ref{sec:sup-sweep}.

\begin{table}[!t]
    \centering
    \caption{Locked MDF-ASMC control parameters used in all numerical tests.}
    \label{sup:control params}
    \begin{tabular}{ll}
    \toprule
        \textbf{Control Parameter} & \textbf{Value}\\
    \midrule
        \multicolumn{2}{l}{\textit{Visibility Funnel --- FoV-Adaptive Cone Clamp}} \\
        $\boldsymbol{\rho}_{\text{fov},0}$      & $[145,105]^\top$~px \\
        $\boldsymbol{\rho}_{\text{fov},\infty}$ & $[40,40]^\top$~px \\
        $\ell_\text{fov}$                       & $0.1$~s$^{-1}$ \\
        $\theta_\text{cap}$                     & $60^\circ$ \\
    \midrule
        \multicolumn{2}{l}{\textit{Normalized-Error PID --- Desired Optic Flow}} \\
        $K_{rp}$ & $\text{diag}([9.0,9.0])$ \\
        $K_{ri}$ & $\text{diag}([0.1,0.1])$ \\
        $K_{rd}$ & $\text{diag}([1.4375,1.4375])$ \\
    \midrule
        \multicolumn{2}{l}{\textit{Performance Function --- Optic Flow}} \\
        $\Xi_2$                     & $\text{diag}([0.2,0.2,0.2])$ \\
        $\boldsymbol{p}_{2_0}$      & $[25.0,25.0,4.0]^\top$ \\
        $\boldsymbol{p}_{2_\infty}$ & $[2.5,2.5,1.5]^\top$ \\
    \midrule
        \multicolumn{2}{l}{\textit{ASMC --- Optic Flow Regulation}} \\
        $\mathcal{X}$           & $\text{diag}([0.05,0.05,0.025])$ \\
        $\Gamma$                & $\text{diag}([0.4375,0.5,0.75])$ \\
        $P$                     & $\text{diag}([1.5,1.5,5.0])$ \\
        $N$                     & $\text{diag}([0.02,0.02,0.05])$ \\
        $\boldsymbol{\kappa}_0$ & $[0.125,0.125,0.25]^\top$ \\
        $E$                     & $\text{diag}([1.0,1.0,1.0])$ \\
        $\epsilon_S$            & $0.05$ \\
    \midrule
        \multicolumn{2}{l}{\textit{Geometric SO(3) Attitude Control}} \\
        $k_R$      & $\text{diag}([1.5,1.5,0.5])$ \\
        $k_\Omega$ & $\text{diag}([0.3,0.3,0.1])$ \\
    \midrule
        \multicolumn{2}{l}{\textit{Adaptive Yaw SMC}} \\
        $\chi_\text{a},\,\gamma_\text{a}$  & $0.5,\,0.5$ \\
        $n_\text{a},\,p_\text{a}$          & $1.0,\,2$ \\
        $\kappa_{\text{a}_0},\,E_\text{a}$ & $2.0,\,3.0$ \\
    \bottomrule
    \end{tabular}
\end{table}

\clearpage
\subsection{Deep-Sweep Gain Analysis}
\label{sec:sup-sweep}

In order to characterise the sensitivity of the locked gain set (Table~\ref{sup:control params}) to parameter perturbations, a cross-trajectory deep sweep was carried out in which each of the $30$ scalar MDF-ASMC parameters (from the outer visibility-funnel gains to the inner geometric SO($3$) and adaptive yaw SMC gains) was independently perturbed by $\{\times 0.5,\times 0.75,\times 1.25,\times 1.5\}$ of its nominal value, and every resulting candidate gain was replayed on all five trajectories for all five initial conditions. This amounts to $30\times 4\times 5\times 5 = 3000$ deterministic simulations (600 unique parameter--multiplier--trajectory combinations, each evaluated over 5 ICs). A candidate is \emph{viable} if it preserves the $5/5$ land rate on every trajectory, and \emph{Pareto-admissible} if it additionally improves either the aggregate landing time $\bar{t}_\text{f}$ or the aggregate worst-case horizontal error $xy_\text{e}^{\max}$ without worsening the other.

The baseline aggregate metrics at the locked gain set are $\bar{t}_\text{f}=18.20~\text{s}$ and $xy_\text{e}^{\max}=0.072~\text{m}$ (mean across the five trajectories). The $120$ parameter--multiplier combinations partition into four classes, summarized in Table~\ref{sup:deep sweep table}.

\begin{table}[!t]
\centering
\caption{Deep-sweep parameter classification across the five trajectories (120 parameter--multiplier combinations, each evaluated on 5 ICs per trajectory).}
\label{sup:deep sweep table}
\begin{tabular}{lcc}
\toprule
\textbf{Class} & \textbf{Count} & \textbf{Fraction} \\
\midrule
Fragile (breaks $5/5$ on $\ge 1$ trajectory) & $53$  & $44\%$ \\
Pareto-admissible (viable $+$ strict improvement) & $42$ & $35\%$ \\
Inert (viable, no measurable change)           & $15$  & $13\%$ \\
Pareto trade (viable, improves one, worsens other) & $10$ & $8\%$ \\
\bottomrule
\end{tabular}
\end{table}

The dominant result is that $44\%$ of all perturbations break the $5/5$ land rate on at least one trajectory, confirming that the majority of the gain set is load-bearing. The $42$ Pareto-admissible perturbations indicate modest headroom for further tuning; however, the locked gains already achieve $25/25$ successful landings with worst-case $xy_\text{e}=7.2~\text{cm}$ (Table~V of the main paper) and are therefore retained as a conservative, validated design point.

The three non-Pareto classes exhibit the following structure.

\emph{Fragile parameters.} Perturbations of $K_{rp}$, $\boldsymbol{p}_{2_0}$, $\boldsymbol{p}_{2_\infty}$, $k_R$, and $\chi_\text{a}$ break the $5/5$ land rate on at least one trajectory --- typically on the lateral-offset initial conditions IC$3$--IC$5$, where the controller's chase authority is most stressed. These parameters are load-bearing at their locked values.

\emph{Inert parameters.} $n_\text{a}$, $p_\text{a}$, $\kappa_{\text{a}_0}$, $E_\text{a}$, $\Omega_\text{a}$, and $\gamma_\text{a}$ produce no measurable change in any reported metric at $\pm 50\%$ perturbation, indicating that the adaptive yaw SMC leakage tail is overparameterized for this class of problems and that the heading-reference integrator is driven almost entirely by the projection error.

\emph{Pareto trades.} A small number of perturbations trade landing time against horizontal precision and thus cannot unilaterally improve both metrics simultaneously.

\clearpage
\subsection{Per-Controller Failure Analysis in the Comparative Benchmark}
\label{sec:sup-comparison}

This section expands on the failure modes of the four baselines in the comparison study of Section~IV-E of the main paper. Table~VI of the main paper reports the quantitative outcomes; the analysis below identifies the structural reason each baseline fails to simultaneously land \emph{precisely} ($xy_\text{e}\le 0.10~\text{m}$) and \emph{softly} ($\|\boldsymbol{v}_\text{rel}\|\le 0.20~\text{m/s}$) under the shared disturbance model of Section~IV-B of the main paper.

For every baseline, the gains reported in the main paper are a best-effort retune \emph{within the framework proposed by the original authors}: the paper-authoritative values of~\cite{lin2022,zhang2026,chen2025,cho2022} each destabilise the closed loop on the adversarial initial condition $\text{IC}_2=[2,2,-5]~\text{m}$ under our disturbance model (crashing within one to three control steps; a per-baseline breakdown of these failures is archived with the implementation code). The retunes therefore isolate a \emph{structural} limitation of each framework rather than a tuning shortcoming.

\subsubsection{Performance-Constrained PBVS \cite{lin2022}}

The scheme of~\cite{lin2022} regulates the position and velocity errors through a single prescribed-performance envelope on each axis,
\begin{equation*}
    \rho_{p_k}(t) = (\rho_{p_k}(0)-\rho_{p_k}^\infty)e^{-l_{p_k}t} + \rho_{p_k}^\infty,
\end{equation*}
with the transformed velocity error $\xi_{v_k}=e_{v_k}/\rho_{v_k}(t)\in(-1,1)$ fed into a barrier function. Widening the $z$-axis envelope to $\rho_{v_\infty}(3)=0.15~\text{m/s}$ absorbs the heave disturbance of~\cite{lin2022} ($A_z=0.2~\text{m}$, $\omega_z=0.5~\text{rad/s}$) used on Cases~2 and~5 (ship-deck), preventing $\xi_{v}(3)$ from saturating on the vertical oscillation. The horizontal channels remain structurally incompatible with the shared disturbance model: the steady-state velocity error induced by the $\pm 0.2~\text{m/s}$ mean wind plus OU turbulence already occupies a substantial fraction of any $\rho_{v_\infty}^{xy}$ compatible with $z$-descent authority, and the additional tracking error from a moving target pushes $\|\xi_v^{xy}\|$ close enough to the barrier that the lateral thrust demand saturates the $60~\text{N}$ cone limit on a subset of noise realisations (observed at $T\approx 49~\text{N}$ on the Linear case). Since the barrier enforces $|\xi_v|<1$ as a hard constraint, there is no gain setting that simultaneously preserves descent authority and admits the worst-case wind-plus-target velocity error; the scheme therefore produces horizontal errors in the $0.71$--$27.3~\text{m}$ range across all five cases, with catastrophic divergence on Case~2.

\subsubsection{PBVS with Adaptive Extended Disturbance Observer \cite{zhang2026}}

The backstepping law of~\cite{zhang2026} realises a second-order closed loop with natural frequency $\omega_n=\sqrt{K_{c_1}K_{c_3}+K_{c_2}}$ and damping $\zeta=(mK_{c_1}+K_{c_3})/(2m\omega_n)$. This baseline is the strongest of the four under our disturbance model, landing on all five trajectories with horizontal error in the $0.08$--$0.59~\text{m}$ band. The structural limitation lies in the decoupling of the horizontal and vertical time scales: the paper-authoritative $\omega_n^z$ drives the UAV through the $5~\text{m}$ descent before the horizontal loop has closed its chase error on moving targets, so landing the adversarial IC inside the $40~\text{s}$ budget forces either a fast $\omega_n^z$ (hard touchdown) or a slow $\omega_n^z$ (horizontal error grows). Our retune ($K_{c_1}^{zz}=0.03$, $K_{c_2}^{zz}=0.15$, $K_{c_3}^{zz}=0.5$, $\omega_n^z\approx 0.4~\text{rad/s}$) gentles the $z$ channel aggressively and still yields terminal relative speeds in the $0.6$--$1.5~\text{m/s}$ range --- well above the $0.20~\text{m/s}$ soft-touchdown threshold. The framework has no mechanism to couple descent-rate regulation to image-plane error closure, so the horizontal--vertical time-scale trade-off cannot be removed by gain tuning.

\subsubsection{IBVS with Adaptive Observer \cite{chen2025}}

The force law of~\cite{chen2025} takes the form $\boldsymbol{f} = -k_r(\boldsymbol{r}-\boldsymbol{s}+\boldsymbol{\varepsilon}) + \text{(observer terms)}$, with $\boldsymbol{r}$ an image-moment vector constructed from the current feature set and $\dot{\boldsymbol{r}}$ required by the depth-adaptive update law. The derivation of~\cite{chen2025} assumes $\dot{\boldsymbol{r}}$ is available analytically from the target motion model; no measurement filter is prescribed. In any vision-only implementation, $\dot{\boldsymbol{r}}$ must be obtained by finite-difference of the noisy feature vector, and the resulting noise energy is amplified by the observer gain and by the $1/z$ scaling implicit in $\boldsymbol{r}$ at low altitude. Under the depth-dependent pixel noise model ($\sigma_{\text{px}}\uparrow$ as $z\downarrow$), the finite-difference error on a non-trivial fraction of random seeds grows without bound as the altitude approaches the landing target, driving the commanded force past the $60~\text{N}$ thrust limit. The scheme successfully lands on Case~1 ($xy_\text{e}=0.08~\text{m}$) but stalls above the landing altitude on all four moving cases: $z_\text{f}=0.39$--$1.24~\text{m}$, with $xy_\text{e}=0.98$--$1.66~\text{m}$, worst on Case~5 ($z_\text{f}=1.24~\text{m}$) where the combined deck motion and target tracking demand amplifies the finite-difference error. Because the stall is a structural consequence of the missing $\dot{\boldsymbol{r}}$ filter rather than a gain-tuning issue, no choice of $k_r$, $k_1$, $k_2$ restores the $5/5$ land rate: lowering $k_r$ (from the paper value of $8$ down to $1.0$) only delays the divergence until lower altitudes without eliminating it.

\subsubsection{Feed-Forward IBVS \cite{cho2022}}

The scheme of~\cite{cho2022} commands a translational velocity $\boldsymbol{v}_d = -\lambda\hat{L}_e^+\boldsymbol{e} + \boldsymbol{v}_{\text{ff}}$ and gates the vertical channel by an adaptive sigmoid altitude gain $a_{d_z}(c) = 1 - \sigma(k_\sigma c)$, where $c$ is the centroid norm of the instantaneous optic-flow vector. In our vision-only harness the velocity feed-forward $\boldsymbol{v}_{\text{ff}}$ is set consistently with the IBVS velocity command. The framework lands on only two of the five cases within the $40~\text{s}$ budget (Cases~2 and~5), with horizontal errors in the $0.26$--$0.43~\text{m}$ range on all trajectories. It cannot simultaneously satisfy the $0.10~\text{m}$ precise and $0.20~\text{m/s}$ soft criteria on any moving trajectory, because the adaptive altitude gain couples the feature-centroid \emph{norm} to the vertical descent command: the paper-authoritative $k_\sigma=0.002$ keeps $a_{d_z}\approx 0.5$ throughout the approach, halving the available descent rate and leaving the UAV $0.4$--$0.5~\text{m}$ above the landing plane when the $40~\text{s}$ budget expires on Cases~1, 3, and~4. Increasing $k_\sigma$ --- e.g.\ to $k_\sigma=0.1$ for fast saturation of $a_{d_z}$ --- breaks moving targets on a different failure mode: the feature centroid jitters under pixel noise, so $a_{d_z}$ jitters and the $z$-command follows, producing the $1$--$3~\text{m}$ lateral errors we observed in that retune. The compromise $k_\sigma=0.02$ used for the reported numbers accepts the terminal stall as the least-bad failure; no setting removes it, because the framework has no separation between precision regulation at altitude and descent regulation near touchdown.

\subsubsection{Comparison Summary Bar Chart}

\begin{figure*}[!t]
\centering
    \includegraphics[width=0.82\textwidth]{Figures/generated/comparison_summary.pdf}
    \caption{Side-by-side summary of landing time, final horizontal error, and UAV-target relative speed at touchdown across all five cases and five controllers. Dotted lines mark the $40~\text{s}$ landing budget, the $0.10~\text{m}$ \emph{precise} horizontal-error bound, and the $0.20~\text{m/s}$ \emph{soft} relative-speed bound. The proposed MDF-ASMC is the only controller that clears all three thresholds simultaneously on every moving trajectory.}
    \label{fig:sup-comp-summary}
\end{figure*}

\clearpage
\subsection{Summary of Numerical Observations}
\label{sec:sup-summary}
The numerical study summarized in Section~IV of the main paper supports the following conclusions:
\begin{enumerate}
    \item The locked MDF-ASMC achieves $25/25$ successful landings across five trajectories and five initial conditions under depth-dependent pixel noise, mean wind plus OU turbulence, pixel outliers, $\pm5\%$ parametric uncertainty, ground effect, computational delay, and actuator constraints, with mean landing time $\le 18.4~\text{s}$ and worst-case horizontal error $\le 7.3~\text{cm}$ --- including ship-deck roll/pitch oscillation and the heave profile from~\cite{lin2022} on Cases~2 and~5. Every touchdown is simultaneously precise ($xy_\text{e}\le 0.10~\text{m}$) and soft ($\|\boldsymbol{v}_\text{rel}\|\le 0.20~\text{m/s}$).
    \item A cross-trajectory deep sweep of $30$ gains $\times$ $4$ perturbations $\times$ $5$ trajectories $\times$ $5$ ICs ($2400$ deterministic runs) partitions the MDF-ASMC parameters into Pareto-admissible, fragile, inert, and Pareto-trade classes. The locked gain set is a justified design point validated by this systematic exploration of the parameter space.
    \item Across the symmetric speed-envelope sweep $\lambda\in[0.6,1.4]$ on all four moving trajectories, every run is simultaneously \emph{precise} ($xy_\text{e}\le 0.10~\text{m}$) and \emph{soft} ($\|\boldsymbol{v}_\text{rel}\|\le 0.20~\text{m/s}$) under the full robustness model, establishing a $\pm 40\%$ band of guaranteed soft-precise touchdown around every nominal trajectory; at the eventual envelope boundary the failure mode is always attitude-cone saturation rather than the adaptive law.
    \item Under identical disturbance conditions, the proposed MDF-ASMC dominates four representative state-of-the-art controllers~\cite{lin2022,zhang2026,chen2025,cho2022} in both landing time and horizontal touchdown precision on all four moving trajectories.
\end{enumerate}

\bibliography{IEEEabrv, bibliography}

\end{document}
